%Root main.tex
%---------------------------------------------------------------------------
%付録A　付録
%---------------------------------------------------------------------------

%\lstnewenvironment{mylisting}[1][]
%    {\lstset{
%        frame=single,
%        numbersep=10pt,
%        tabsize=2,
%        #1
%    }
%}{}
\begin{mylisting}[language=c,caption=hold\_function2,label=appC_hold_func1]
function y1  = fcn(u1,cnt)
\%永続変数の定義
persistent y1\_buf;

\%初期化
if isempty(y1\_buf)
    y1\_buf = 0;
end

\%演算
if(cnt == 0)
    y1 = u1;
    y1\_buf = u1;
else
    y1 = y1\_buf;
end

\end{mylisting}

リスト\ref{appC_hold_func1}はキャリアカウンタであるcntが0の時のみ値を更新し，cntが0以外の場合は，0の時点での値を出力する内容となっている．hold\_function\_thetaも同様となっている．hold\_function\_thetaのプログラムを以下に示す．


\begin{mylisting}[language=c,caption=hold\_function\_theta,label=appC_hold_func2]
function out  = fcn(in1,cnt,sector)
\%永続変数の定義
persistent in1\_buf sector\_buf;

\%初期化
if isempty(in1\_buf)
    in1\_buf = 0;
end

if isempty(sector\_buf)
    sector\_buf = 0;
end

\%演算
if(cnt == 0)
    sector\_buf = sector;
    in1\_buf = in1;
    out = in1;
else
    if(sector == sector\_buf)
        out = in1;
        in1\_buf = in1;
    else
        out = in1\_buf;
    end
end

\end{mylisting}
